\documentclass[12pt,a4paper]{ctexart}
\usepackage{geometry}
\geometry{left=2.5cm,right=2.5cm,top=2.5cm,bottom=2.5cm}
\usepackage{graphicx}
\usepackage{amsmath,amssymb}
\usepackage{float}
\usepackage{subfigure}
\usepackage{booktabs}
\usepackage{hyperref}

\title{透明塑料饮料瓶内残留液体识别的深度学习模型研究}
\author{参赛队伍}
\date{\today}

\begin{document}
\maketitle

\begin{abstract}
本文针对透明塑料饮料瓶内残留液体的识别与定位问题，建立了基于深度学习和计算机视觉的自动化分析模型。首先，我们构建了包含多样化场景的数据集；其次，利用 PyTorch 搭建了用于裁剪图像区域分类的卷积神经网络，实现了高精度的液体有无研判；最后，进一步推广至端到端的 YOLO 目标检测框架，实现了瓶体的自动定位与液体含量的三分类（少、中、多）预测。

\textbf{关键词：} 图像分类，目标检测，YOLO，深度学习，液位估算
\end{abstract}

\newpage
\tableofcontents
\newpage

\section{问题重述与背景分析}
\subsection{问题背景}
\subsection{问题重述}
针对题意，我们需要完成：
1. 构建不少于 100 张、多样性丰富的透明瓶数据集并标注。
2. 针对确定边界框区域，对其进行特征提取并分类是否有液体。
3. 建立通用的模型，一键实现区域识别（Bounding Box 定位）与液体容量判定。

\section{模型假设与符号说明}
\subsection{模型假设}
\begin{itemize}
    \item 假设人工拍摄环境下的瓶身形变不足以阻碍卷积神经网络的特征提取。
    \item 假设所有液体含量的级别可以依据像素面积与其占总瓶身的比例划分为定性区间。
\end{itemize}
\subsection{符号说明}
在此列出如 $B_{box}$, $L_p$, $IoU$ 等相关参数和函数符号。

\section{任务一：多场景数据集的构建}
\subsection{图像采集策略}
说明是如何采集覆盖不同姿态（竖立、横放、倾斜等）、不同液体颜色透明度的图片。
\subsection{数据标注与增强 (Data Augmentation)}
介绍使用标注工具划分 Bounding Box，并阐述引入几何变换、亮度增强对于提高模型泛化能力的作用。

\section{任务二：基于 CNN 的残留液体二分类模型}
\subsection{模型选择与网络架构}
对比 ResNet、MobileNet 等，阐述最终选择方案的原因及网络层设计，写出关键的前向传播数学表达式等。
\subsection{损失函数与优化算法}
阐述所使用的二分类交叉熵损失函数（Binary Cross-Entropy Loss）。
$$ L = -\frac{1}{N}\sum_{i=1}^{N}[y_i \log(\hat{y}_i) + (1-y_i)\log(1-\hat{y}_i)] $$
\subsection{实验结果与泛化能力评价}
给出在测试集上的 Accuracy, Precision, Recall 和 F1-Score。绘制混淆矩阵（Confusion Matrix）。

\section{任务三：基于 YOLO 的自动定位与容量等级评估}
\subsection{目标检测框架设计}
说明如何将问题转化为回归加多分类问题，YOLO 的锚框机制（Anchor）与损失函数的构建（含有 Bbox Loss, Class Loss 和 DFL Loss）。
\subsection{多级别融合标注方式}
通过合并“无液体，少，中，多”四类直接让 YOLO 对边界框输出多标签，或先验定位瓶身再对截取图像进行容量比例评估。
\subsection{泛化与鲁棒性测试}
展示 YOLO 模型在未知测试集上的定位 IoU 分数及内容量分类的正确率，提供可视化效果图（可在此处插入带框效果图）。

\section{模型评价与改进方向}
\subsection{模型的优点}
\subsection{模型的缺点与局限}

\section{结论}
总结各项任务的达成度与实际在自动分拣和垃圾分类中的工业落地应用潜力。

\begin{thebibliography}{99}
\bibitem{ref1} Redmon, J., \& Farhadi, A. (2018). YOLOv3: An Incremental Improvement. \textit{arXiv preprint arXiv:1804.02767}.
\bibitem{ref2} He, K., Zhang, X., Ren, S., \& Sun, J. (2016). Deep residual learning for image recognition. \textit{CVPR}.
\end{thebibliography}

\end{document}